
\subsection{A Timestamped Asynchronous Tool Benchmark for On-Device Agents}
Our system contributions in Design 1--2 target a specific but increasingly common deployment: on-device agents that interleave LLM planning with multiple tool calls whose latencies are heterogeneous, stochastic, and often longer than a single prefill pass.
To evaluate such systems, we need a benchmark that exposes the \emph{temporal dimension} of agent execution, i.e., when actions are issued and when observations arrive, rather than only the semantic correctness of the final answer.
However, existing tool-use datasets and benchmarks are not designed for this setting.
ToolBench~\cite{qin2023toolllm} provides large-scale trajectories grounded in real REST APIs and tool outputs, but it is primarily organized as synchronous, order-preserving tool interaction.
AsyncTool~\cite{shi2026asynctool} formalizes asynchronous multi-task tool calling with delayed tool responses and evaluates agents under time-layered execution, but it does not aim to preserve \emph{real} API outputs for every call.
DTA-LLaMA~\cite{zhu2025divideThenAggregate} focuses on parallel tool invocation by converting linear tool traces into a DAG and training agents to divide-then-aggregate, but it does not provide fine-grained wall-clock timestamps for out-of-order arrivals.

\textbf{Key idea.}
We propose a new benchmark and dataset construction principle: \emph{retain ToolBench as the source of tool-call semantics (tool name, arguments, and real outputs), and augment it with an explicit asynchronous timeline (timestamps, parallelism, and out-of-order returns) using an AsyncTool-style temporal protocol, optionally with DTA-style DAG structure for dependencies.}
This directly matches the on-device agent scenario studied in this paper, where overall responsiveness is determined by how well the system overlaps LLM compute with tool I/O.

\subsubsection{Task Definition and Data Schema}
Each sample is a time-stamped execution trace consisting of alternating \emph{planning segments} and \emph{tool-call events}.
Let $t$ denote wall-clock time (relative to the start of the user request).
We represent one episode as an ordered event stream $\mathcal{E} = \{e_1,\ldots,e_N\}$ where each event is one of:
(i) user input, (ii) model planning output, (iii) tool invocation, (iv) tool return, (v) final model response.

Concretely, a sample contains:
(1) \textbf{User request:} text $U$ with timestamp $t_U$.
(2) \textbf{Planning segments:} a sequence $\{(P_r, t^{\text{start}}_{P_r}, t^{\text{end}}_{P_r})\}_{r=1}^{R}$, where $P_r$ is the model-generated plan/update text for round $r$.
(3) \textbf{Asynchronous tool calls:} a set of calls $\{c_k\}_{k=1}^{K}$.
Each call $c_k$ records tool name, serialized arguments, and two timestamps $(t^{\text{start}}_k, t^{\text{end}}_k)$ together with the tool output $O_k$ (from ToolBench) and an explicit \emph{return order} induced by sorting $t^{\text{end}}_k$.
(4) \textbf{Final answer:} text $A$ with timestamp $t_A$.

This schema supports the interaction pattern we target: $U \rightarrow P_1 \rightarrow \{c_k\} \rightarrow \text{(out-of-order } O_k\text{)} \rightarrow P_2 \rightarrow \cdots \rightarrow A$.
Importantly, we store both a machine-readable event list (for metric computation) and a linearized transcript (for prompting), where tool returns are inserted into the context strictly by their observed arrival times.

\textbf{Oracle-plan track (aligned with this paper).}
To isolate system-level scheduling and KV-management effects, we include an \emph{oracle-plan} variant in which the planning segments and tool-call arguments are treated as correct by construction.
This matches our focus: given correct calls and real tool outputs, can the agent \emph{execute} efficiently and integrate out-of-order observations without breaking correctness?

\subsubsection{Construction Pipeline}
We build the benchmark by extending ToolBench trajectories~\cite{qin2023toolllm} with an asynchronous timeline.
The pipeline is:
(1) \textbf{Trajectory selection and normalization.} Select multi-tool samples and normalize tool I/O into a canonical serialized format.
(2) \textbf{Dependency extraction (optional DAG).} Convert the tool sequence into a DAG whose nodes are tool calls and whose edges represent explicit data dependencies.
We follow the divide-then-aggregate perspective of DTA-LLaMA~\cite{zhu2025divideThenAggregate}: calls without dependencies are eligible for parallel invocation.
(3) \textbf{Latency injection and timestamping.} Assign each call a duration $\Delta_k$ and set $t^{\text{end}}_k = t^{\text{start}}_k + \Delta_k$.
Durations can be drawn from fitted distributions (e.g., per-category long-tail) or from device-specific measurements.
This step instantiates the temporal protocol emphasized by AsyncTool~\cite{shi2026asynctool} while keeping ToolBench outputs intact.
(4) \textbf{Transcript synthesis.} Emit the interleaved, arrival-ordered transcript with planning segments and tool observations, producing the exact context stream that the on-device system would see.

\subsubsection{Metrics for On-Device Asynchronous Agents}
Our benchmark is designed to evaluate both \emph{correctness} and \emph{time efficiency}.
In addition to standard answer-level correctness (e.g., ToolBench/ToolEval-style success)~\cite{qin2023toolllm}, we introduce temporal metrics:
(1) \textbf{Episode makespan:} $T_{\text{end}} = t_A - t_U$.
(2) \textbf{Time-to-first-token (TTFT):} the time from $t_U$ to the start of final decoding.
(3) \textbf{Idle-time ratio:} the fraction of time the device spends waiting for tool returns without useful LLM compute, estimated from execution logs.
(4) \textbf{Scheduling overhead:} number of interaction rounds and redundant tool calls (to discourage brute-force switching).

These metrics directly reflect the objective of Design 2: by running incremental prefill upon each arrival, LAP reduces idle time and TTFT while preserving correctness under out-of-order tool observations.